\documentclass{article}

\usepackage{corl_2026} % Use this for the initial submission.
% \usepackage[final]{corl_2026} % Uncomment for the camera-ready ``final'' version.
% \usepackage[preprint]{corl_2026} % Uncomment for pre-prints (e.g., arxiv); This is like ``final'', but will remove the CORL footnote.

\title{Fault-Tolerant Control of Robotic Manipulators via Deep Reinforcement Learning}

% The \author macro works with any number of authors. There are two
% commands used to separate the names and addresses of multiple
% authors: \And and \AND.
%
% Using \And between authors leaves it to LaTeX to determine where to
% break the lines. Using \AND forces a line break at that point. So,
% if LaTeX puts 3 of 4 authors names on the first line, and the last
% on the second line, try using \AND instead of \And before the third
% author name.

% NOTE: authors will be visible only in the camera-ready and preprint versions (i.e., when using the option 'final' or 'preprint'). 
% 	For the initial submission the authors will be anonymized.

\author{
  Jane E.~Doe\\
  Department of Electrical Engineering and Computer Sciences\\
  University of California Berkeley 
  United States\\
  \texttt{janedoe@berkeley.edu} \\
  %% examples of more authors
  %% \And
  %% Coauthor \\
  %% Affiliation \\
  %% Address \\
  %% \texttt{email} \\
  %% \AND
  %% Coauthor \\
  %% Affiliation \\
  %% Address \\
  %% \texttt{email} \\
  %% \And
  %% Coauthor \\
  %% Affiliation \\
  %% Address \\
  %% \texttt{email} \\
  %% \And
  %% Coauthor \\
  %% Affiliation \\
  %% Address \\
  %% \texttt{email} \\
}


\begin{document}
\maketitle

%===============================================================================

\begin{abstract}
	Robotic manipulators operating in industrial and safety-critical environments are vulnerable to joint
	failures that can abruptly degrade performance, yet existing fault-tolerant strategies rely on accurate
	dynamic models or predefined recovery procedures that do not generalize across diverse failure modes.
	We present COMPENSATE, a deep reinforcement learning framework in which specialized teacher policies
	trained on individual failure configurations are distilled into a unified student policy that detects
	and adapts to joint failures in real time, without any explicit fault signal. By incorporating
	permanent lockouts, degraded stiffness, and simultaneous multi-joint faults into the training
	distribution, COMPENSATE learns compensation strategies that transfer to unseen failure types. We
	evaluate on reach and lift tasks in IsaacLab with cross-simulator validation in mjlab and zero-shot
	deployment on a Franka Research~3 arm. Against standard PPO, a diffusion policy baseline, and a
	fault-aware inverse kinematics controller, COMPENSATE achieves superior task completion rates and
	adaptation speed while maintaining smooth operational transitions under diverse joint failure conditions.
\end{abstract}

% Two or three meaningful keywords should be added here
\keywords{Fault-Tolerant Control, Deep Reinforcement Learning, Robotic Manipulation, Sim-to-Real Transfer, Teacher-Student Distillation}

%===============================================================================

\section{Introduction}

	Robotic manipulators are a cornerstone of modern industrial applications, performing repetitive,
	high-precision tasks with speed and reliability~\citep{ali2018vision, ghodsian2023mobile}. Recent
	advances in artificial intelligence~(AI) have expanded the capabilities of robotic systems, enabling
	them to learn, adapt, and make independent decisions~\citep{sun2021artificial, liu2021deep,
	li2024manipllm}. Yet despite these advances, a fundamental vulnerability persists: mechanical
	malfunctions (arising from joint wear, actuator fatigue, or unexpected environmental loading) can
	abruptly degrade performance or cause complete task failure~\citep{291826, sciavicco2012modelling}.
	As manipulators are deployed in safety-critical domains such as healthcare, manufacturing, and field
	robotics, the ability to maintain task performance through joint failures has become essential, while
	most deployed systems still rely on predefined, fault-unaware control
	policies~\citep{deloitte2025smart, allaboutai2025report}.

	Traditional fault management strategies address this vulnerability through hardware redundancy and
	fault diagnosis and isolation~(FDI). Redundancy-based approaches add actuators or parallel kinematic
	chains so that a single joint failure does not halt operation~\citep{291826, anand2012fault,
	mokhtari2022fault}, but incur substantial cost, weight, and design complexity. FDI techniques
	(spanning signal-processing anomaly detection~\citep{kimevaluation, ompusunggu2014winding}, rule-based
	expert systems~\citep{li2023open}, and model-based deviation monitors~\citep{shuai2022research,
	sabry2019fault}) can identify that a fault has occurred, yet they do not synthesize an adapted control
	policy in response; a separate, hand-crafted recovery layer is still required.

	Machine learning has begun to close this gap. Predictive maintenance and anomaly detection models can
	anticipate failures before they occur or flag deviations during operation~\citep{lian2022anomaly,
	eski2011fault}, and recent work demonstrates that learning directly from observed robot failures
	(including recovery in long-horizon manipulation tasks) is feasible~\citep{huang2025fail2progress}.
	Generative approaches such as counterfactual failure synthesis can further produce diverse fault
	rollouts from successful demonstrations, reaching 81.3\% correction accuracy~\citep{wang2025learning}.
	Nevertheless, these methods presuppose either a corpus of failure demonstrations or an explicit
	fault-detection signal as input; when a joint degrades silently mid-task and no labeled failure data
	are available, they cannot generate an online, adaptive control response.

	The robotics community has increasingly turned to reinforcement learning~(RL) to overcome this
	limitation. For legged locomotion, contrastive forward-prediction RL enables agents to detect joint
	damage from proprioceptive signals alone and immediately reconfigure gait without any explicit fault
	label~\citep{fu2025contrastive}. For visuomotor policies, latent-space RL can steer a pretrained
	diffusion policy toward successful outcomes with minimal additional
	samples~\citep{wagenmaker2025steering}, and sim-to-real transfer via massive parallel sampling with
	active domain exploration has demonstrated 42--63\% performance gains over strong
	baselines~\citep{sobanbabu2025sampling}. These results confirm that RL-based adaptation is a
	principled direction for fault-tolerant robotics, yet they only address locomotion. The specific
	challenge of reactive joint-failure compensation for \emph{robot manipulator} reaching tasks in real
	hardware deployments remains underexplored.

	To address this gap, we present \textbf{COMPENSATE}, a deep reinforcement learning framework for
	fault-tolerant control of robotic manipulators. COMPENSATE employs a \emph{teacher-student}
	architecture: specialized teacher policies are first trained on individual failure configurations,
	then distilled into a unified student policy that adapts in real time without any explicit fault
	signal~\citep{pham2024adaptive}. Our contributions are threefold. \emph{(i)}~A teacher-student
	distillation scheme that generalizes across \emph{permanent lockouts}, \emph{degraded stiffness},
	and \emph{simultaneous multi-joint faults}, covering both seen and unseen failure modes at test time.
	\emph{(ii)}~Sim-to-real transfer demonstrated by zero-shot deployment on a Franka Research~3 arm,
	with cross-simulator validation in mjlab confirming policy robustness beyond the IsaacLab training
	environment. \emph{(iii)}~Comprehensive benchmarking against standard PPO, a diffusion policy
	baseline, and a fault-aware inverse kinematics controller, showing superior task completion and
	adaptation speed across reach and lift tasks.

%===============================================================================

\section{Related Work}
\label{sec:related}

	\textbf{Classical fault-tolerant control.} Fault-tolerant control for robotic manipulators has long
	relied on model-based methods. Early strategies applied passive and active hardware redundancy so that
	joint failures would not halt operation, while fault diagnosis and isolation~(FDI) detected deviations
	from known system models~\citep{blanke, 291826, sciavicco2012modelling}. More recent designs embed
	backstepping or sliding-mode controllers that maintain tracking under actuator
	faults~\citep{anand2012fault, mokhtari2022fault, act11120353}. These methods are analytically rigorous
	but require accurate dynamics models and hand-crafted recovery rules, limiting their applicability to
	anticipated, well-characterized fault modes.

	\textbf{Learning-based fault detection and diagnosis.} Data-driven methods have substantially improved
	fault detection speed and accuracy. Neural networks, deep belief networks, and power-signature analysis
	can flag incipient anomalies before hard failures occur~\citep{eski2011fault, sabry2019fault,
	shuai2022research, raouf2022mechanical}. Inverse reinforcement learning has been applied to
	simultaneously detect anomalies and propose corrective actions in autonomous
	systems~\citep{lian2022anomaly}, and deep-learning diagnosis has been extended to multi-fault
	industrial robots~\citep{adam2023multiple}. Detection alone, however, does not yield an adapted
	control policy; a separate recovery layer is still required, and these methods are unreliable when
	failures evolve gradually with no labeled fault data available.

	\textbf{Reinforcement learning for fault-tolerant robots.} RL has emerged as a compelling framework
	for learning adaptive behaviors without explicit failure models. In legged locomotion, contrastive
	forward-prediction RL enables quadrupeds to detect joint damage from proprioceptive signals alone and
	immediately reconfigure gait without any fault label~\citep{fu2025contrastive}. Learning from observed
	failure trajectories has also shown promise for manipulation recovery~\citep{huang2025fail2progress,
	wang2025learning}, and RL has been compared favorably against model predictive control for
	fault-tolerant continuous control~\citep{ahmed2018comparison}. For manipulators specifically,
	\citet{pham2024adaptive} formulated joint-failure compensation as a partially observable Markov
	decision process and demonstrated high task success, but addressed only single-joint, simulation-only
	scenarios, leaving multi-fault and real-hardware settings open.

	\textbf{Teacher-student distillation and sim-to-real transfer.} Teacher-student RL has become the
	standard paradigm for deploying policies that must operate under partial observability at test time.
	A privileged teacher trained with full state access (e.g., exact fault vectors or terrain properties)
	is distilled into a student that operates on deployable observations alone~\citep{kumar2021rma}.
	Latent-space RL can further steer a pre-trained policy toward correct behavior with minimal additional
	samples~\citep{wagenmaker2025steering}, and active-exploration system identification during sim-to-real
	transfer yields 42--63\% performance gains over standard domain
	randomization~\citep{sobanbabu2025sampling}. COMPENSATE applies this paradigm to manipulator fault
	tolerance: fault-specialized teacher policies are distilled into a single student that compensates for
	joint failures zero-shot on physical hardware.


%===============================================================================

\section{Method}
\label{sec:method}

	\textbf{COMPENSATE} is a two-phase teacher-student framework for fault-tolerant manipulator
	control. In Phase~1, specialized teacher policies are trained via PPO on individual
	joint-failure configurations. In Phase~2, a single student policy is distilled from the
	teacher ensemble via behavioral cloning, without access to any explicit fault signal at
	deployment.

	\subsection{Problem Formulation}
	\label{subsec:problem}

		We model joint-failure-resilient manipulation as a Partially Observable Markov Decision
		Process~(POMDP) $\langle\mathcal{S},\mathcal{O},\mathcal{A},T,R,\gamma\rangle$. The full
		state $s \in \mathcal{S}$ includes joint positions, joint velocities, and the current
		failure configuration. The policy receives only a 32-dimensional proprioceptive
		observation $o \in \mathcal{O}$: relative joint positions $q \in \mathbb{R}^{9}$, relative
		joint velocities $\dot{q} \in \mathbb{R}^{9}$, the previous action
		$a_{t-1} \in \mathbb{R}^{7}$, and a goal end-effector pose $g \in \mathbb{R}^{7}$
		(3D position and unit quaternion). No fault indicator is included in $o$; the agent must
		infer the failure state purely from observable kinematics.

		Actions $a \in \mathcal{A} \subset \mathbb{R}^{7}$ are relative joint-position increments
		with scale $\alpha = 0.03$~rad. We consider $|\mathcal{F}|=8$ failure configurations:
		seven single-joint permanent lockouts $\{f_0,\ldots,f_6\}$ and a healthy baseline
		$f_\emptyset$. Under lockout $f_i$, joint $i$ holds its last commanded position regardless
		of the network output:
		\begin{equation}
		    q_t^{\mathrm{cmd}, i} =
		        \begin{cases}
		            q_{t-1}^{\mathrm{cmd}, i} + \alpha\,a_t^i & \text{joint } i \text{ healthy,}\\
		            q_{t-1}^{\mathrm{cmd}, i}                  & \text{joint } i \text{ locked.}
		        \end{cases}
		    \label{eq:action}
		\end{equation}

	\subsection{Teacher Policy Training}
	\label{subsec:teachers}

		For each configuration $c \in \mathcal{F}$, we train a specialist teacher~$\pi^c$ using
		Proximal Policy Optimization~\citep{schulman2017proximal}. Each teacher is a two-layer MLP
		with hidden dimensions $[128,\,64]$ and ELU activations. Training runs for $3{,}500$
		iterations with 24 rollout steps per parallel environment, an adaptive learning rate
		(target KL\,$=0.01$, initial $\eta = 10^{-4}$), clip ratio $\epsilon=0.2$, and entropy
		coefficient $\beta=0.001$.

		The per-step reward combines end-effector tracking and motion regularization:
		\begin{equation}
		    r_t = w_\mathrm{pos}\,\|p_\mathrm{ee} - p_g\|_2
		          + w_\mathrm{fine}\,r_\mathrm{fine}(p_\mathrm{ee},\,p_g)
		          + w_\mathrm{ori}\,r_\mathrm{ori}(\mathbf{R}_\mathrm{ee},\,\mathbf{R}_g)
		          + w_\mathrm{reg}\,r_\mathrm{reg}(a_t,\,\dot{q}_t),
		    \label{eq:reward}
		\end{equation}
		where $p_\mathrm{ee}$, $p_g$ are end-effector and goal positions,
		$\mathbf{R}_\mathrm{ee}$, $\mathbf{R}_g$ are their rotation matrices, and
		$r_\mathrm{fine}$ is a kernel-smoothed fine-positioning bonus. The fixed weights are
		$w_\mathrm{pos}=-2.0$ and $w_\mathrm{ori}=-1.0$. The fine-grained weight
		$w_\mathrm{fine}$ and action-rate and joint-velocity regularizers are each annealed from
		zero to their target values over the first $5{,}000$ training steps, encouraging the
		agent to first learn to reach the goal before optimizing for motion smoothness.

		To ensure all commanded targets are reachable under the active joint constraints, goal
		poses are generated by sampling random joint configurations, computing forward kinematics,
		and discarding poses outside the reachable workspace of the locked joint.

	\subsection{Fault Simulation and Curriculum}
	\label{subsec:curriculum}

		Joint failures are simulated at the action level: at each episode reset, one configuration
		$c$ is drawn uniformly from the current curriculum set $\mathcal{C}_k \subseteq \mathcal{F}$.
		During teacher training, an incremental curriculum initializes $\mathcal{C}_k$ with only
		the teacher’s target failure mode and expands by one additional mode every
		$\lfloor 1/|\mathcal{F}|\rfloor$ fraction of total training steps, so the policy first
		masters a single failure type before encountering the full distribution. This staged
		exposure stabilizes early-training dynamics. During student distillation~(Phase~2),
		$\mathcal{C}_k = \mathcal{F}$ throughout with no curriculum gating, since all teachers
		are already fully trained.

	\subsection{Multi-Teacher Distillation}
	\label{subsec:distillation}

		After all teachers are trained, we distill them into a single student policy~$\pi^S$ via
		behavioral cloning. The student is a wider MLP $[128,\,128]$ with ELU activations and
		receives the same 32-dimensional observation as the teachers; no failure indicator is
		present in its input.

		During distillation rollouts, each parallel environment is independently assigned a random
		configuration $c$ at reset, and the corresponding teacher~$\pi^c$ provides the behavioral
		target for that environment. Given a minibatch
		$\{(o_k,\,a_k^c)\}_{k=1}^{B}$, the student minimizes:
		\begin{equation}
		    \mathcal{L}_\mathrm{BC} =
		        \frac{1}{B}\sum_{k=1}^{B}
		        \bigl\|\pi^S(o_k) - a_k^c\bigr\|_2^{2}.
		    \label{eq:bc}
		\end{equation}
		Training uses the Adam optimizer ($\eta = 5\times10^{-4}$), gradient accumulation over
		15 steps, gradient clipping at norm~1.0, and 8 learning epochs per rollout. Because the
		failure index is withheld from~$\pi^S$, the student is forced to recover a compensatory
		strategy purely from proprioceptive dynamics, enabling generalization to failure modes not
		encountered during teacher training.

	\subsection{Simulation Environments and Sim-to-Real Transfer}
	\label{subsec:sim2real}

		All policies are trained in \textbf{IsaacLab} on a Franka Research~3~(FR3) model. The
		reach task commands random end-effector poses resampled every 6~s, with success defined
		as reaching within 0.04~m position error and 0.2~rad orientation error. Additive uniform
		observation noise ($\pm 0.01$~rad) is applied to joint positions, velocities, and action
		feedback at every step to encourage robustness.

		Cross-simulator validation is performed in \textbf{mjlab}~(MuJoCo~Warp), which mirrors
		the IsaacLab observation and action specification at a matched 30~Hz control rate
		(240~Hz physics, decimation~8). Checkpoints trained in IsaacLab transfer to mjlab
		without fine-tuning, confirming that learned compensation behaviors are not artifacts of
		the Isaac~Sim physics engine. Zero-shot deployment on a physical Franka Research~3 arm
		completes the sim-to-real evaluation.



\section{Experiments -- 1-2 Pages}
\subsection{Implementation Details}
We conducted our simulation experiments using IsaacLab, which offers high-fidelity physics simulation for accurate modeling of robot manipulators and their interactions with the environment. 
Within this environment, we trained both teacher and student policies from scratch using proximal policy optimization (PPO). The training process incorporated task-specific reward functions as well as curricula designed to progressively increase task difficulty, ensuring that the learned policies were both robust and adaptable.

Following policy training in IsaacLab, we performed cross-simulator evaluation using mjlab. This sim-to-sim transfer step allowed us to examine how the trained policies generalize under different physics engines. Since MuJoCo employs dynamics models that differ from those in IsaacLab, successful performance across both platforms indicates that the policies would transfer   to real hardware. Finally, we validated our approach in hardware on two separate manipulators, the Franka Research 3 (FR3) and the Franka Emika Panda. These real-world trials facilitated a comprehensive assessment of zero-shot policy deployment across both simulation and physical domains.

\subsection{Tasks and Evaluation Protocol}

To evaluate the COMPENSATE framework, we defined a suite of manipulation tasks designed to capture core operational capabilities and resilience under failure conditions. Specifically, we considered four representative tasks: Reach, Spline, Lift, and Stack. Collectively, these tasks span the essential skill set of articulated manipulators, from precise target acquisition to trajectory execution, object lifting, and sequential object placement.

To investigate fault tolerance, we introduced two distinct classes of simulated joint malfunctions. The first was a permanent joint failure, modeled by constraining the affected joint to its last commanded position and enforcing zero velocity, thereby emulating a joint that ceases operation entirely. The second was a degraded performance fault, implemented by artificially increasing joint stiffness, resulting in diminished accuracy and compliance while preserving partial operability. These failure models were injected during task execution, forcing the learned policies to adapt dynamically to unexpected hardware limitations.

The trained policies were evaluated under both seen tasks, which were included during training, and unseen tasks, which were held out to assess generalization. This dual evaluation protocol enabled us to systematically investigate the robustness and adaptability of the learned control strategies when confronted with novel scenarios beyond their training distribution. For all tasks, the primary evaluation criteria consisted of task completion rate and completion time, both of which serve as direct indicators of efficiency and reliability. In addition, for the Spline task, we incorporated a trajectory-tracking metric to quantify positional accuracy with respect to the reference path. This allowed us to capture not only binary success or failure but also the degree of precision in executing continuous, dynamically constrained movements.

\subsection{Results}

To benchmark the proposed method, we compared its performance against several baselines: a standard PPO policy trained without joint failures, a diffusion policy, and an inverse kinematics (IK)-based controller. Following the training of the teacher policies, we generated datasets encompassing all modeled joint failure scenarios and subsequently distilled this data into a diffusion policy, thereby providing a generative baseline for fault adaptation. In addition, we developed a fault-tolerant IK system in which the index and position of the failed joint were explicitly provided to the kinematics solver. This allowed the solver to reconfigure the manipulator and compute feasible joint configurations under constrained degrees of freedom, serving as a classical control baseline for handling hardware faults.


%===============================================================================

\section{Experimental Results}
\label{sec:result}

	Nam dui ligula, fringilla a, euismod sodales, sollicitudin vel, wisi. Morbi auctor lorem non justo.
	Nam lacus libero, pretium at, lobortis vitae, ultricies et, tellus. Donec aliquet, tortor sed accumsan
	bibendum, erat ligula aliquet magna, vitae ornare odio metus a mi. Morbi ac orci et nisl hendrerit
	mollis. 

	Suspendisse ut massa. Cras nec ante. Pellentesque a nulla. Cum sociis natoque penatibus
	et magnis dis parturient montes, nascetur ridiculus mus. Aliquam tincidunt urna. Nulla ullamcorper
	vestibulum turpis. Pellentesque cursus luctus mauris.
	
	Nulla malesuada porttitor diam. Donec felis erat, congue non, volutpat at, tincidunt tristique, libero.
	Vivamus viverra fermentum felis. Donec nonummy pellentesque ante. Phasellus adipiscing semper 	elit. 
	Proin fermentum massa ac quam. Sed diam turpis, molestie vitae, placerat a, molestie nec, leo.
	Maecenas lacinia. Nam ipsum ligula, eleifend at, accumsan nec, suscipit a, ipsum. Morbi blandit
	ligula feugiat magna. Nunc eleifend consequat lorem. Sed lacinia nulla vitae enim. Pellentesque tin-
	cidunt purus vel magna. Integer non enim. Praesent euismod nunc eu purus. Donec bibendum quam
	in tellus. Nullam cursus pulvinar lectus. Donec et mi. Nam vulputate metus eu enim. Vestibulum
	pellentesque felis eu massa.
	Quisque ullamcorper placerat ipsum. Cras nibh. Morbi vel justo vitae lacus tincidunt ultrices. 
	
	Lorem
	ipsum dolor sit amet, consectetuer adipiscing elit. In hac habitasse platea dictumst. Integer tempus
	convallis augue. Etiam facilisis. Nunc elementum fermentum wisi. Aenean placerat. Ut imperdiet,
	enim sed gravida sollicitudin, felis odio placerat quam, ac pulvinar elit purus eget enim. Nunc vitae
	tortor. Proin tempus nibh sit amet nisl. Vivamus quis tortor vitae risus porta vehicula.

%===============================================================================

\section{Conclusion}
\label{sec:conclusion}

	Nam dui ligula, fringilla a, euismod sodales, sollicitudin vel, wisi. Morbi auctor lorem non justo.
	Nam lacus libero, pretium at, lobortis vitae, ultricies et, tellus. Donec aliquet, tortor sed accumsan
	bibendum, erat ligula aliquet magna, vitae ornare odio metus a mi. Morbi ac orci et nisl hendrerit
	mollis. Suspendisse ut massa. Cras nec ante. Pellentesque a nulla. Cum sociis natoque penatibus
	et magnis dis parturient montes, nascetur ridiculus mus. Aliquam tincidunt urna. Nulla ullamcorper
	vestibulum turpis. Pellentesque cursus luctus mauris.
	Nulla malesuada porttitor diam. Donec felis erat, congue non, volutpat at, tincidunt tristique, libero.
	Vivamus viverra fermentum felis. Donec nonummy pellentesque ante. Phasellus adipiscing semper
	elit. Proin fermentum massa ac quam. Sed diam turpis, molestie vitae, placerat a, molestie nec, leo.
	Maecenas lacinia. Nam ipsum ligula, eleifend at, accumsan nec, suscipit a, ipsum. Morbi blandit
	ligula feugiat magna. Nunc eleifend consequat lorem. Sed lacinia nulla vitae enim. Pellentesque tin-
	cidunt purus vel magna. Integer non enim. Praesent euismod nunc eu purus. Donec bibendum quam
	in tellus. Nullam cursus pulvinar lectus. Donec et mi. Nam vulputate metus eu enim. Vestibulum
	pellentesque felis eu massa.
	Quisque ullamcorper placerat ipsum. Cras nibh. Morbi vel justo vitae lacus tincidunt ultrices. Lorem
	ipsum dolor sit amet, consectetuer adipiscing elit. In hac habitasse platea dictumst. Integer tempus
	convallis augue. Etiam facilisis. Nunc elementum fermentum wisi. Aenean placerat. Ut imperdiet,
	enim sed gravida sollicitudin, felis odio placerat quam, ac pulvinar elit purus eget enim. Nunc vitae
	tortor. Proin tempus nibh sit amet nisl. Vivamus quis tortor vitae risus porta vehicula.

%===============================================================================

\clearpage
% The acknowledgments are automatically included only in the final and preprint versions of the paper.
\acknowledgments{If a paper is accepted, the final camera-ready version will (and probably should) include acknowledgments. All acknowledgments go at the end of the paper, including thanks to reviewers who gave useful comments, to colleagues who contributed to the ideas, and to funding agencies and corporate sponsors that provided financial support.}

%===============================================================================

% no \bibliographystyle is required, since the corl style is automatically used.
\bibliography{example}  % .bib

\end{document}
